% §14.1 The Optimality Question
%
% Monograph register: frame the chapter against the chapter-8 floor and
% Bloom gap, pose the constructive question, sketch the BHF answer. Cuts
% the textbook roadmap and second-person framing. Cross-refs
% thm:entropy-floor, sec:entropy-floor, prop:bloom-overhead, sec:bloom-gap,
% def:bloom-filter, sec:bhf-construction, def:bhf, sec:space-analysis,
% sec:fpr-analysis. DEFINED: sec:optimality-question.

\section{The Optimality Question}
\label{sec:optimality-question}

Chapter~8 closed Part~II with a contrast it could not resolve. The
entropy floor (\Cref{thm:entropy-floor}) certified that some FNR-zero
representation could in principle reach $-\log_2 \fprate$ bits per
element, but it exhibited none; the standard Bloom filter, the only
concrete FNR-zero scheme on the table, sat a fixed multiplicative
factor above that floor. The bound was diagnostic, not constructive.
This chapter supplies the construction: an explicit FNR-zero
approximate set that meets the floor up to integer rounding and a
sublinear correction, closing the gap the Bloom filter leaves open.

\subsection*{The gap}

\Cref{thm:entropy-floor} of \Cref{sec:entropy-floor} states that any
FNR-zero approximate-set representation with false-positive rate at
most $\fprate$ requires
\[
  L / n \;\ge\; -\log_2 \fprate \;+\; o(1)
\]
bits per element in the asymptotic regime $N \gg n$, $\fprate$ fixed.
The argument is a counting bound: $2^L$ representations must cover
$\binom{N}{n}$ possible target sets, and each representation describes
at most $\binom{|A_R|}{n}$ true sets inside its accepted set, forcing
$L$ to grow linearly in $n \log_2(1/\fprate)$.

\Cref{prop:bloom-overhead} of \Cref{sec:bloom-gap} computes the Bloom
filter's budget from \Cref{def:bloom-filter}: at the optimal hash count
$k^* = (m/n)\ln 2$, the per-element cost is
\[
  m / n \;=\; \log_2(1/\fprate) \cdot \log_2 e,
\]
or $\log_2 e \approx 1.4427$ times the floor, a constant independent of
$\fprate$. The overhead traces to a structural choice: $k$ independent
hashes write into a flat shared bit array, and the analysis maximizes
per-bit entropy by tuning the array to half occupancy. Half-occupancy
is the wrong target. Bits should pay for the discrimination the
structure performs, not for the encoding's internal balance.

\subsection*{The constructive question}

\Cref{sec:bloom-gap} isolated three Bloom design choices, $k$
independent hashes, a flat shared bit array, and a conjunction query,
and observed that closing the gap requires abandoning at least one.
Which one, and what replaces it, the remark left open. A construction
that answers it needs four properties: FNR-zero by construction, so it
stays on the perfect-filter axis with the Bloom filter; an FPR provably
at most $\fprate$; a per-element budget matching $-\log_2 \fprate$ up to
a sublinear correction; and a query running in time comparable to
Bloom's. The first two are correctness, the third is the optimality the
chapter promises, the fourth is what makes the construction deployable
rather than an existence proof. The floor theorem guarantees that some
encoding meets all four asymptotically; the question is whether an
explicit one does, and what it looks like.

\subsection*{The answer in sketch}

The Bernoulli Hash Function (BHF) abandons the flat shared bit array.
In its place it uses a \emph{perfect hash} $\phi : A \to [n]$ that
bijects the designated set $A$ onto $[n] = \{1, \ldots, n\}$, and a
\emph{fingerprint table} $F : [n] \to \{0, 1\}^r$ storing
$F[\phi(a)] = g(a)$ for each $a \in A$, where $g : \Universe \to
\{0, 1\}^r$ is a secondary hash and $r = \lceil -\log_2 \fprate \rceil$
is the fingerprint width. A query for $x$ computes the slot $\phi(x)$,
reads the stored fingerprint $F[\phi(x)]$, and returns $\top$ exactly
when it equals $g(x)$.

For $x \in A$ the perfect hash sends $x$ to its own slot, where the
build pass stored $g(x)$; the comparison always succeeds, so no false
negatives occur. For $x \notin A$ the slot holds $g(a)$ for whichever
$a \in A$ shares the slot, and $g(a) = g(x)$ with probability $2^{-r}$
under the uniformity of $g$; with $r = \lceil -\log_2 \fprate \rceil$
the false-positive rate is at most $\fprate$ by construction. The space
cost is $n r$ bits for the table plus a sublinear perfect-hash term, so
the per-element budget is
\[
  r \;+\; o(1)
  \;=\; -\log_2 \fprate \;+\; O(1),
\]
matching the floor up to the integer rounding of $r$. The construction
trades Bloom's flat array for a perfect-hash-addressed table, and the
per-element cost drops from $\log_2 e \cdot \log_2(1/\fprate)$ to
$\log_2(1/\fprate)$ plus rounding. \Cref{sec:bhf-construction} gives the
construction in full, \Cref{sec:space-analysis} the bit budget and the
optimality theorem, and \Cref{sec:fpr-analysis} the false-positive
bound.
